



\section{Introduction}

% Multi-armed bandits (MABs) have garnered significant attention as they provide a well-defined framework and techniques for investigating the trade off between experimentation (learning) and rewards (earning) and in sequential decision-making problems. In MAB problems, a decision-maker sequentially selects actions from a given set and observes corresponding uncertain rewards. The MAB setting also extends to scenarios where the decision rewards can be personalized by the presence of features or covariates, leading to contextual bandits, as showcased in a large number of recent applications \citep{langford2008theepoch,li2010contextual,tewari2017ads,zhou2020survey,villar2015multi,bastani2020online,cohen2020feature}.

Multi-armed bandits (MABs) provide a framework for studying the exploration-exploitation trade-off in sequential decision-making, where a decision-maker selects actions and observes uncertain rewards. This extends to contextual bandits with features or covariates, as shown in numerous applications \citep{langford2008theepoch,li2010contextual,tewari2017ads,zhou2020survey,villar2015multi,bastani2020online,cohen2020feature}. This paper focuses on a well-studied class of models that captures both MABs and contextual bandits as special cases while being amenable to theoretical analysis: the stochastic linear bandit (LB) problem. In this model, the problem parameter $\param$ represents an unknown vector in $\RR^d$, while the actions, also vectors in $\RR^d$, yield noisy rewards with a mean equal to the inner product of $\param$ and the chosen action. The objective of a policy is to maximize the cumulative reward based on the observed data up to the decision time. The policy's performance is measured by the cumulative regret, which quantifies the difference between the total expected rewards achieved by the policy and the maximum achievable expected reward. 

Achieving this objective necessitates striking a balance between exploration and exploitation. In the context of LB, this entails selecting actions that aid in estimating the true parameter $\param$ accurately while obtaining optimal rewards. Various algorithms based on the \emph{optimism principle} have been developed to address this challenge, wherein the optimal action is chosen based on the upper confidence bound (UCB) \citep{lai1985asymptotically, auer2002using, dani2008stochastic, rusmevichientong2010linearly}. Another popular strategy is Thompson sampling (TS), a Bayesian heuristic introduced by \cite{thompson1933likelihood} that employs randomization to select actions according to the posterior distribution of reward functions. Additionally, the Greedy policy that selects the myopically best action is shown to be effective in contextual bandits \citep{kannan2018smoothed,raghavan2018externalities,hao2020adaptive, bastani2021mostly}.

% In the linear bandit setting, two types of regret are considered. The Bayesian regret is applicable when the parameter $\param$ is treated as a random variable with a prior distribution. In this case, the regret is averaged over three sources of randomness: the observation noise, the randomized algorithm, and the random parameter $\param$. Intuitively, the Bayesian regret measures the expected performance of the algorithm over different realizations of the parameter $\param$. \cite{russo2014learning} and \cite{dong2018information} establish an $\tilde{\cO}(d \sqrt{T})$ upper bound for the Bayesian regret of the Thompson Sampling (TS) heuristic, referred to as LinTS, which matches the minimax optimal bound shown by \cite{dani2008stochastic}. Here, $\tilde{\cO}$ refers to the asymptotic order, up to polylogarithmic factors.

% On the other hand, the frequentist regret assumes a fixed, unknown parameter $\param$, and the expectation is taken only with respect to the randomness of the noise and the algorithm. In this frequentist setting, the optimism-based algorithm proposed by \cite{abbasi2011improved}, known as Optimism in the Face of Uncertainty Linear Bandit Algorithm (OFUL), achieves a $\tilde{\cO}(d\sqrt{T})$ frequentist regret bound, which matches the minimax optimal Bayesian bound. However, for a frequentist variant of LinTS, referred to as TS-Freq, which modifies the posterior distribution by increasing its variance, \cite{agrawal2013thompson} and \cite{abeille2017linear} provide a $\tilde{\cO}(d\sqrt{dT})$ upper bound for the frequentist regret. This regret bound falls short of the optimal rate by a factor of $\sqrt{d}$. Recent work by \citet{hamidi2020frequentist} confirms that this modification is necessary, and thus, the frequentist regret of LinTS cannot be improved. For the Greedy algorithm in linear bandit problems, no general theoretical guarantee exists \citep{lattimore2017end}, and hence both LinTS and Greedy might perform suboptimally.



In the linear bandit setting, two regret types are considered. Bayesian regret treats parameter $\param$ as a random variable with a prior distribution, averaging regret over noise, algorithm randomness, and parameter randomness, measuring expected performance across parameter realizations. \cite{russo2014learning} and \cite{dong2018information} establish an $\tilde{\cO}(d \sqrt{T})$ upper bound for the  Bayesian regret of the Thompson Sampling (TS) heuristic, referred to as LinTS, matching the minimax optimal bound by \cite{dani2008stochastic}. Here, $\tilde{\cO}$ denotes asymptotic order up to polylogarithmic factors. Frequentist regret assumes fixed unknown parameter $\param$, averaging only over noise and algorithm randomness. The OFUL algorithm \citep{abbasi2011improved} achieves an optimal $\tilde{\cO}(d\sqrt{T})$ frequentist regret bound. However, TS-Freq, a frequentist LinTS variant with inflated posterior variance, only achieves $\tilde{\cO}(d\sqrt{dT})$ \citep{agrawal2013thompson,abeille2017linear}, suboptimal by factor $\sqrt{d}$. \citet{hamidi2020frequentist} confirms this inflation is necessary and LinTS's frequentist regret cannot be improved. The Greedy algorithm lacks theoretical guarantees for linear bandit problems \citep{lattimore2017end}, suggesting both LinTS and Greedy may perform suboptimally.


Despite the theoretical gaps, LinTS demonstrates strong empirical performance \citep{russo2018tutorial}, suggesting posterior distribution inflation may be unnecessary in most scenarios. Similarly, the Greedy algorithm performs well in typical cases \citep{bietti2021contextual}. While optimism-based algorithms are computationally expensive (generally NP-hard \citep{dani2008stochastic,russo2014learning,agrawal2019recent}), LinTS and Greedy maintain computational efficiency. This disparity between theoretical, computational, and empirical performance prompts two questions: Can we identify problematic instances for LinTS and Greedy in a data-driven way and apply "course-correction" to ensure competitive frequentist regret bounds? Can this be achieved while preserving their empirical performance and computational efficiency?
In this paper, we provide positive answers to both questions. Specifically, we make the following \emph{contributions}.

\noindent 1. We develop a real-time geometric analysis technique for the $d$-dimensional confidence ellipsoid surrounding $\param$. This method is crucial for maximizing the use of historical data, advancing beyond methods that capture only limited information from the confidence ellipsoid, such as a single numerical value. Consequently, this facilitates a more precise ``course-correction''.

\noindent 2.  We introduce a comprehensive family of algorithms, termed \emph{\POFUL{}} (encompassing OFUL, LinTS, TS-Freq, and Greedy as specific instances), and derive a general, data-driven frequentist regret bound for them. This bound is efficiently computable using data observed from previous decision epochs.

\noindent 3.   We introduce course-corrected variants of LinTS and Greedy that achieve minimax optimal frequentist regret. These adaptations maintain most of the desirable characteristics of the original algorithms.
%    \end{enumerate}

%The distinguishing feature of our study from the existing literature is the geometric analysis of the $d$-dimensional confidence ellipsoid. 	This is in contrast to techniques outlined in previous studies that may inadvertently discard valuable information as a consequence of their focus on summary statistics.


%The distinguishing feature of our study from the existing literature is the geometric analysis of the $d$-dimensional confidence ellipsoid. Specifically, we conduct an analysis considering the size, shape, position, and orientation of the confidence ellipsoid during the execution of bandit algorithms. This enables real-time characterization of the set of potentially optimal actions, leading to the establishment of a data-driven regret bound.


\subsection{Other Related Literature}

Our work is closely related to three main research streams: methodological foundations of linear bandits, bandit algorithms utilizing spectral properties, and data-driven exploration techniques. While these works share some similarities with our approach, we highlight the key differences and the unique aspects of our methodology.

From a methodological perspective, our regret analysis builds upon the foundations laid by \cite{abbasi2011improved}, \cite{agrawal2013thompson}, and \cite{abeille2017linear}. However, a key distinguishing factor is that our approach does not rely on optimistic samples, which is a departure from previous methods. This means that the algorithms we study do not always choose actions that are expected to perform better than the true optimal action. By allowing non-optimistic samples, we avoid the need to inflate the posterior distribution, a requirement in the works of \cite{agrawal2013thompson} and \cite{abeille2017linear}.

Our use of spectral information in bandit algorithms bears some resemblance to the study of \emph{Spectral Bandits} \citep{valko2014spectral, kocak2014Spectral, kocak2020spectral, kocak2020Best}. These works represent arm rewards as smooth functions on a graph, leveraging low-rank structures to improve algorithmic performance and obtain regret guarantees independent of the number of actions. In contrast, our approach exploits the spectral properties of the action covariance matrix, which is distinct from graph spectral analysis. Moreover, our research tackles the broader context of stochastic linear bandits without assuming any low-rank structure.

Our work also shares conceptual similarities with research on exploration strategies \citep{russo2016information, kirschner2018heteroscedastic} and data-driven exploration reduction \citep{bastani2021mostly, pacchiano2020regret, hamidi2020frequentist, hamidi2020general}. However, our methodology and data utilization differ significantly. For instance, \cite{bastani2021mostly} focuses on the minimum eigenvalue of the covariance matrix, a single-parameter summary of the observed data, while \cite{hamidi2020general} uses information from one-dimensional reward confidence intervals. The work of \cite{hamidi2020frequentist} is more closely related to ours, as it employs spectral information to improve the performance of Thompson Sampling in linear bandits. They use a single summary statistic called the \emph{thinness coefficient} to decide whether to inflate the posterior. In contrast, our approach leverages the full geometric details of the $d$-dimensional confidence ellipsoid, harnessing richer geometric information.

% \subsection{Organization}

% The structure of this paper is as follows: Section \ref{sec:setting} introduces the notations and problem formulation. In Section \ref{sec:POFUL}-\ref{sec:regret}, we present the \POFUL{} family of algorithms and analyze their frequentist regret. 
% Sections \ref{sec:data-driven}-\ref{sec:meta} detail our data-driven approach and meta-algorithm for course-correction. Section  \ref{sec:simulation} discusses simulations using synthetic and real-world datasets. The paper concludes in Section \ref{sec:conclusion} with a summary of findings. All technical proofs, further discussions, and supplementary materials are provided in the appendices.

